% Generated draft: requires mathematical and source review.
\documentclass[11pt]{article}
\usepackage[margin=25mm]{geometry}
\usepackage{fontspec}
\setmainfont{Latin Modern Roman}
\usepackage{amsmath,amssymb,mathtools,unicode-math}
\setmathfont{Latin Modern Math}
\usepackage{enumitem,needspace,xurl,hyperref}
\setlength{\parindent}{0pt}
\setlength{\parskip}{5pt}
\setlength{\emergencystretch}{3em}
\newcommand{\N}{\mathbb{N}}
\newcommand{\Z}{\mathbb{Z}}
\newcommand{\Q}{\mathbb{Q}}
\newcommand{\R}{\mathbb{R}}
\newcommand{\C}{\mathbb{C}}
\newcommand{\F}{\mathbb{F}}
\newcommand{\A}{\mathbb{A}}
\newcommand{\PP}{\mathbb{P}}
\newcommand{\E}{\mathbb{E}}
\newcommand{\eps}{\varepsilon}
\newcommand{\dd}{\,\mathrm{d}}
\begin{document}
\section*{UnsolvedMath ID: 30002100}
\textbf{Source problem code:} catalogue entry
\subsection{Definitions and mathematical statement}
For an algebraically closed field \(k\) and an integer \(g\geq 2\), an unmarked Deligne--Mumford stable curve of genus \(g\) over \(k\) is a connected, reduced, projective curve \(C\) with only nodes as singularities, arithmetic genus \(\dim_k H^1(C,\mathcal O_C)=g\), and ample dualizing sheaf \(\omega_C\). Write \(\overline{\mathcal M}_{g,k}\) for the moduli stack of such curves and \(\overline M_{g,k}\) for its coarse moduli space. These definitions explain the standard unmarked notation; the source does not specify the base field, its characteristic, or the intended quantification over genera.

A minimal model program for a specified canonical or log canonical divisor seeks a sequence of divisorial contractions and flips that are negative for that divisor, ending in a minimal model, where the resulting divisor is nef, or a Mori fibre space. Here nef means having nonnegative intersection with every complete integral curve. In the usual variety formulation, a Mori fibre space has a positive-dimensional general fibre, relative Picard number one, and the negative of the relevant canonical or log canonical divisor relatively ample. Applying this description requires specifying the underlying variety or stack, the divisor, and the applicable hypotheses.

\textbf{Source question.} Can the minimal model program be run on \(\overline{\mathcal M}_g\), the moduli space of Deligne--Mumford stable curves?

The supplied background connects this question with the Hassett--Keel program and the construction and interpretation of log canonical models and flips. A corresponding possible interpretation concerns log canonical divisors \(K+\alpha\Delta\), where \(K\) is the canonical divisor on the chosen moduli object and \(\Delta\) is its boundary of singular stable curves. However, the source does not specify this divisor family, a parameter range, or whether the intended target is the ordinary canonical minimal model program, the full log canonical variation, or modular descriptions of its models. It also calls \(\overline{\mathcal M}_g\) a space without resolving the stack versus coarse-space distinction. Consequently, the displayed question is retained as a programmatic question: a unique fully quantified formulation requires these choices to be resolved.
\subsection{Short English statement}
The question asks whether one can carry out the minimal model program for the moduli of stable curves. The supplied background points toward the Hassett--Keel program, but the intended genera, base field, stack or coarse space, divisor, and completion criterion remain unspecified. The supplied \emph{partially solved} label is dated curation, not an independently verified conclusion.
\subsection{Sources}
\begin{itemize}
\item Local registry: \texttt{data/unsolvedmath/problems.json}, permanent UnsolvedMath ID \(30002100\), exact source problem code \texttt{catalogue entry}, title \emph{The Minimal Model Program for Stable-Curve Moduli}; collection: Oberwolfach Reports Open Problems. No individual proposer is supplied. The recorded proposed year is \(2012\).
\item Supplied source citation: \emph{Classical Algebraic Geometry (2013)}, Oberwolfach Reports, DOI \texttt{10.4171/owr/2012/30}, pp.~1845--1893. \url{https://doi.org/10.4171/owr/2012/30}. The citation's title year \(2013\) differs from the recorded proposed year \(2012\); this discrepancy is preserved without resolution.
\item Supplied literature source: Hassett and Hyeon, first flip. \url{https://annals.math.princeton.edu/2013/177-3/p03}.
\item Supplied literature source: Hyeon, outline of the log MMP. \url{https://arxiv.org/abs/1006.1094}.
\item Supplied literature source: original OWR report. \url{https://doi.org/10.4171/owr/2012/30}.
\item Literature assessment dated exactly \texttt{2026-08-21}: the supplied curation labels the problem partially solved and reports that the Hassett--Keel program has constructed and interpreted initial log canonical models and flips, while the full minimal model program for stable-curve moduli remains incomplete in general. This is supplied dated curation derived from the OWR material, not independently verified here; its time-sensitive status and cited claims require independent verification.
\item The supplied record labels the statement \texttt{corrected\_verified} and reports that its mathematical content was checked against the official Oberwolfach report, with notation, context, and question form normalized without changing the source-level claim. This is a supplied source-verification assertion, not verification performed for this draft.
\end{itemize}
\end{document}
